% GENERATED FILE. Edit canonical lessons, then run npm run generate:books.
% canonical-source-hash: fnv1a64:f3718b0cfca0b76e
% canonical-lessons: GU-C361-ninda, GU-C361-tika, GU-C361-thapko, GU-C361-vagovvu, GU-C361-vinavvu

\chapter{Criticism and Pleading}
\label{ch:gu-criticism-and-pleading}
\hlchaptermodality{361}

\begin{chapteropening}
\textbf{By the end of this chapter:} \emph{I can criticise someone, scold them or plead with them.}

Complete the last lesson of chapter 361: I can criticise someone, scold them or plead with them.
\end{chapteropening}

\section[\texorpdfstring{\gu{નિંદા} (nindā)}{નિંદા (nindā)}]{\gu{નિંદા} (nindā) — speaking ill of someone}

\label{lesson:GU-C361-ninda}



\begin{quote}
Before the new one: say the Gujarati for a betrayal, then the Gujarati for to entrap.
\end{quote}

\subsection*{You'll want to know: \gu{નિંદા}}
\textbf{\gu{નિંદા}} — \emph{nindā} — \textquotedblleft{}speaking ill of someone\textquotedblright{}.

\begin{quote}
\gu{બીજાની} \gu{નિંદા} \gu{કરવી} \gu{સારી} \gu{વાત} \gu{નથી}.
\end{quote}

\begin{quote}
\emph{bījānī nindā karvī sārī vāt nathī.}
\end{quote}

\begin{quote}
It isn't good to speak ill of others.
\end{quote}

A feminine noun.

\begin{cousinweb}
\textbf{\gu{નિંદા}} (\emph{nindā}) is a tatsama, borrowed from Sanskrit \emph{nindā}, \textquotedblleft{}blame, censure, reproach\textquotedblright{} (Monier-Williams).
\end{cousinweb}

\begin{grammarlens}[title={What you've built}]
One more word for giving reasons and opinions.
\end{grammarlens}

\subsection*{Your turn}
\begin{itemize}
\raggedright
  \item \emph{Say it:} \emph{nindā}
  \item \emph{Say it:} \emph{nindā}, once more
  \item \emph{Say it:} \emph{bījānī nindā karvī sārī vāt nathī.} — It isn't good to speak ill of others.
  \item \emph{Recall:} say the Gujarati for a betrayal, then the Gujarati for to entrap, then say \emph{nindā} again
\end{itemize}

\begin{tcolorbox}[breakable,colback=teal!4,colframe=teal!35!black,title={Before you move on}]
What does \gu{નિંદા} mean? (\textquotedblleft{}Speaking ill of someone\textquotedblright{}.) Say it once more.
\end{tcolorbox}
\section[\texorpdfstring{\gu{ટીકા} (ṭīkā)}{ટીકા (ṭīkā)}]{\gu{ટીકા} (ṭīkā) — criticism}

\label{lesson:GU-C361-tika}



\begin{quote}
Before the new one: say the Gujarati for to entrap, then the Gujarati for speaking ill of someone.
\end{quote}

\subsection*{You'll want to know: \gu{ટીકા}}
\textbf{\gu{ટીકા}} — \emph{ṭīkā} — \textquotedblleft{}criticism\textquotedblright{}.

\begin{quote}
\gu{જનતાએ} \gu{સરકારની} \gu{નવી} \gu{નીતિની} \gu{ટીકા} \gu{કરી}.
\end{quote}

\begin{quote}
\emph{jantāe sarkārnī navī nītinī ṭīkā karī.}
\end{quote}

\begin{quote}
The public criticised the government's new policy.
\end{quote}

A feminine noun.

\begin{cousinweb}
\textbf{\gu{ટીકા}} (\emph{ṭīkā}) is a tatsama, borrowed from Sanskrit \emph{ṭīkā}, \textquotedblleft{}a commentary\textquotedblright{} (Monier-Williams); Gujarati has moved it from comment to criticism.
\end{cousinweb}

\begin{grammarlens}[title={What you've built}]
One more word for giving reasons and opinions.
\end{grammarlens}

\subsection*{Your turn}
\begin{itemize}
\raggedright
  \item \emph{Say it:} \emph{ṭīkā}
  \item \emph{Say it:} \emph{ṭīkā}, once more
  \item \emph{Say it:} \emph{jantāe sarkārnī navī nītinī ṭīkā karī.} — The public criticised the government's new policy.
  \item \emph{Recall:} say the Gujarati for to entrap, then the Gujarati for speaking ill of someone, then say \emph{ṭīkā} again
\end{itemize}

\begin{tcolorbox}[breakable,colback=teal!4,colframe=teal!35!black,title={Before you move on}]
What does \gu{ટીકા} mean? (\textquotedblleft{}Criticism\textquotedblright{}.) Say it once more.
\end{tcolorbox}
\section[\texorpdfstring{\gu{ઠપકો} (ṭhapko)}{ઠપકો (ṭhapko)}]{\gu{ઠપકો} (ṭhapko) — a telling-off}

\label{lesson:GU-C361-thapko}



\begin{quote}
Before the new one: say the Gujarati for speaking ill of someone, then the Gujarati for criticism.
\end{quote}

\subsection*{You'll want to know: \gu{ઠપકો}}
\textbf{\gu{ઠપકો}} — \emph{ṭhapko} — \textquotedblleft{}a telling-off\textquotedblright{}.

\begin{quote}
\gu{હું} \gu{મોડો} \gu{આવ્યો} \gu{એટલે} \gu{ગુરુએ} \gu{મને} \gu{ઠપકો} \gu{આપ્યો}.
\end{quote}

\begin{quote}
\emph{hũ moḍo āvyo eṭle gurue mane ṭhapko āpyo.}
\end{quote}

\begin{quote}
I came late, so the teacher told me off.
\end{quote}

A masculine noun.

\begin{cousinweb}
\textbf{\gu{ઠપકો}} (\emph{ṭhapko}) CDIAL lists the Gujarati noun, \textquotedblleft{}a rebuke\textquotedblright{}, under a reconstructed \emph{ṭhapp}, \textquotedblleft{}strike, press\textquotedblright{} (CDIAL 5495).
\end{cousinweb}

\begin{grammarlens}[title={What you've built}]
One more word for giving reasons and opinions.
\end{grammarlens}

\subsection*{Your turn}
\begin{itemize}
\raggedright
  \item \emph{Say it:} \emph{ṭhapko}
  \item \emph{Say it:} \emph{ṭhapko}, once more
  \item \emph{Say it:} \emph{hũ moḍo āvyo eṭle gurue mane ṭhapko āpyo.} — I came late, so the teacher told me off.
  \item \emph{Recall:} say the Gujarati for speaking ill of someone, then the Gujarati for criticism, then say \emph{ṭhapko} again
\end{itemize}

\begin{tcolorbox}[breakable,colback=teal!4,colframe=teal!35!black,title={Before you move on}]
What does \gu{ઠપકો} mean? (\textquotedblleft{}A telling-off\textquotedblright{}.) Say it once more.
\end{tcolorbox}
\section[\texorpdfstring{\gu{વગોવવું} (vagovvũ)}{વગોવવું (vagovvũ)}]{\gu{વગોવવું} (vagovvũ) — to run down (to defame)}

\label{lesson:GU-C361-vagovvu}



\begin{quote}
Before the new one: say the Gujarati for criticism, then the Gujarati for a telling-off.
\end{quote}

\subsection*{You'll want to know: \gu{વગોવવું}}
\textbf{\gu{વગોવવું}} — \emph{vagovvũ} — \textquotedblleft{}to run down (to defame)\textquotedblright{}.

\begin{quote}
\gu{તે} \gu{બધાની} \gu{સામે} \gu{પોતાના} \gu{ભાઈને} \gu{વગોવે} \gu{છે}.
\end{quote}

\begin{quote}
\emph{te badhānī sāme potānā bhāīne vagove che.}
\end{quote}

\begin{quote}
He runs his brother down in front of everyone.
\end{quote}

Past, with \gu{તેણે} (he or she …): \gu{તેણે} \gu{પોતાના} \gu{પડોશીને} \gu{વગોવ્યો}, \emph{teṇe potānā paḍośīne vagovyo}.

\begin{cousinweb}
\textbf{\gu{વગોવવું}} (\emph{vagovvũ}) CDIAL lists the Gujarati verb, \textquotedblleft{}to revile\textquotedblright{}, under a reconstructed \emph{vigrōcati}, \textquotedblleft{}robs\textquotedblright{} (CDIAL 11671): to rob someone of a good name.
\end{cousinweb}

\begin{grammarlens}[title={What you've built}]
One more word for giving reasons and opinions.
\end{grammarlens}

\subsection*{Your turn}
\begin{itemize}
\raggedright
  \item \emph{Say it:} \emph{vagovvũ}
  \item \emph{Say it:} \emph{vagovvũ}, once more
  \item \emph{Say it:} \emph{te badhānī sāme potānā bhāīne vagove che.} — He runs his brother down in front of everyone.
  \item \emph{Recall:} say the Gujarati for criticism, then the Gujarati for a telling-off, then say \emph{vagovvũ} again
\end{itemize}

\begin{tcolorbox}[breakable,colback=teal!4,colframe=teal!35!black,title={Before you move on}]
What does \gu{વગોવવું} mean? (\textquotedblleft{}To run down (to defame)\textquotedblright{}.) Say it once more.
\end{tcolorbox}
\section[\texorpdfstring{\gu{વીનવવું} (vīnavvũ)}{વીનવવું (vīnavvũ)}]{\gu{વીનવવું} (vīnavvũ) — to entreat}

\label{lesson:GU-C361-vinavvu}



\begin{quote}
Before the new one: say the Gujarati for a telling-off, then the Gujarati for to run down.
\end{quote}

\subsection*{You'll want to know: \gu{વીનવવું}}
\textbf{\gu{વીનવવું}} — \emph{vīnavvũ} — \textquotedblleft{}to entreat\textquotedblright{}.

\begin{quote}
\gu{મીરા} \gu{દાક્તરને} \gu{વહેલા} \gu{આવવા} \gu{વીનવે} \gu{છે}.
\end{quote}

\begin{quote}
\emph{Mīrā dāktarne vahelā āvvā vīnve che.}
\end{quote}

\begin{quote}
Mira is pleading with the doctor to come early.
\end{quote}

Past, with \gu{તેણે} (he or she …): \gu{તેણે} \gu{ઘણું} \gu{વીનવ્યું}, \emph{teṇe ghaṇũ vīnavyũ}.

\begin{cousinweb}
\textbf{\gu{વીનવવું}} (\emph{vīnavvũ}) is inherited from Sanskrit \emph{víjñapayati}, \textquotedblleft{}makes known, begs\textquotedblright{}; CDIAL lists the Gujarati verb, \textquotedblleft{}to entreat\textquotedblright{} (CDIAL 11705).
\end{cousinweb}

\begin{grammarlens}[title={What you've built}]
One more word for giving reasons and opinions.
\end{grammarlens}

\subsection*{Your turn}
\begin{itemize}
\raggedright
  \item \emph{Say it:} \emph{vīnavvũ}
  \item \emph{Say it:} \emph{vīnavvũ}, once more
  \item \emph{Say it:} \emph{Mīrā dāktarne vahelā āvvā vīnve che.} — Mira is pleading with the doctor to come early.
  \item \emph{Recall:} say the Gujarati for a telling-off, then the Gujarati for to run down, then say \emph{vīnavvũ} again
\end{itemize}

\begin{tcolorbox}[breakable,colback=teal!4,colframe=teal!35!black,title={Before you move on}]
What does \gu{વીનવવું} mean? (\textquotedblleft{}To entreat\textquotedblright{}.) Say it once more.
\end{tcolorbox}
